\subsection{An integral reverse map for the Fermat-eight decomposition}

\begin{statement}{A.188}{Proved}
{A literal multiplier-sixteen reverse isogeny and a five-level Selmer bound}
Let \(F:J_8\to\mathcal B\) be exactly the isogeny
(A.182.3).  Write \(f_i:C_8\to E_i\),
\(g_j:C_8\to G_j\), and \(h_l:C_8\to D_l\)
for the three, six, and three curve maps in that proof,
in its displayed order.  Put
\(j_l:\mathcal P_l\hookrightarrow\operatorname{Jac}(D_l)\),
and \((c_1,\ldots,c_6)=(2,2,2,1,1,1)\).
The following is a homomorphism of abelian varieties over \(\Q\):
\[
\begin{aligned}
 H:\mathcal B&\longrightarrow J_8,\\
 ((e_i),(z_j),(p_l))&\longmapsto
    \sum_{i=1}^3 f_i^*e_i
       +\sum_{j=1}^6 c_jg_j^*z_j
       +\sum_{l=1}^3 h_l^*j_lp_l.
\end{aligned}
\tag{A.188.1}
\]
Here every star in (A.188.1) denotes pullback of degree-zero
divisor classes along the specified curve map.  It satisfies
the integral identities
\[
                HF=[16]_{J_8},\qquad FH=[16]_{\mathcal B}.
\tag{A.188.2}
\]
Consequently \(\ker F\subset J_8[16]\),
\(\ker H\subset\mathcal B[16]\), and \(\deg F\)
is a power of two.  In the arithmetic situation already proved
in A.185--A.187, the explicit consequences are
\[
\begin{aligned}
 32\Sha(J_8/\Q)[2^\infty]&=0,\\
 \operatorname{im}\bigl(
   \operatorname{Sel}^{2^{n+5}}(J_8/\Q)
       \longrightarrow\operatorname{Sel}^{2^n}(J_8/\Q)\bigr)
   &=\kappa_{2^n}\bigl(J_8(\Q)/2^nJ_8(\Q)\bigr)
       \qquad(n\geq1).
\end{aligned}
\tag{A.188.3}
\]
The transition in the second line is induced by \([32]\)
on the finite torsion modules.  All these identities concern
the actual rational homomorphisms and commute with localization.
\end{statement}

\begin{proof}
\emph{Exact degrees of all twelve curve maps.}
We first determine their generic geometric fibers.  All counts
below take place on the dense open where the displayed
denominators and root arguments are nonzero.  Roots in
characteristic zero are distinct there.

For the elliptic maps, first use the quartic coordinates
\((x,y)\) in (A.182.5).  The fibers of the three maps
to those quartics are respectively specified by
\[
\begin{array}{c|c|c|c}
\text{row}&\text{first choice}&\text{second choice}
       &\text{number of choices}\\ \hline
1&v^4=1/y&u^2=xv^2&4\cdot2=8\\
2&v^2=x&u^4=y/3&2\cdot4=8\\
3&u^2=x&v^4=y/2&2\cdot4=8.
\end{array}
\tag{A.188.4}
\]
The quartic equation verifies \(3u^8-2v^8=1\)
for every one of the listed choices, and the forward
coordinates recover exactly \((x,y)\).  The additional
map \((x,y)\mapsto(ax^2,axy)\) has degree two:
its inverse consists of \(x^2=X/a\),
\(y=Y/(ax)\).  Thus \(\deg f_i=16\) for all three.

For the genus-two maps, use the intermediate coordinates
\((x,y)\) in (A.182.7), before
\((X,Y)=(x^2,xy)\).  Their fibers are
\[
\begin{array}{c|c|c|c}
j&\text{root choice}&\text{remaining coordinate}
    &\text{number of choices}\\ \hline
1&v^4=1/y&u=xv&4\\
2&u^4=y/3&v=x&4\\
3&v^4=y/2&u=x&4\\
4&u^8=(y+1)/6&v=x/u&8\\
5&v^8=1/(y-1)&u=xv^2&8\\
6&u^8=2/(3-y)&v=xu^2&8.
\end{array}
\tag{A.188.5}
\]
Substitution of \(y^2=Ax^8+B\) checks the source
equation and recovers both forward coordinates in each
row.  For instance, row four uses
\(u^8=(y+1)/6\) and \(v^8=(y-1)/4\), since
\(y^2-1=24x^8\).  In row five,
\(y-1=1/v^8\); in row six,
\(3-y=2/u^8\).  The subsequent map has exactly
two inverse choices, \(x^2=X\), \(y=Y/x\).
This proves
\[
                 (\deg g_j)_{j=1}^6=(8,8,8,16,16,16).
\tag{A.188.6}
\]

The three quartic maps in (A.182.8) have fibers
\[
\begin{array}{c|c|c|c}
l&\text{first choice}&\text{remaining choice or coordinate}
      &\text{number of choices}\\ \hline
1&u^4=-2z&v^2=w/u&4\cdot2=8\\
2&v^4=3z&u^2=w/v&4\cdot2=8\\
3&v^8=z/w^4&u=wv^3&8.
\end{array}
\tag{A.188.7}
\]
The equation \(w^4=Az^3+Bz\) verifies the source
equation and all forward coordinates.  Hence every
\(h_l\) has degree eight.  These are generic-fiber
counts for the specified morphisms, and so also prove
their degrees over \(\Q\).

\emph{Divisor pullback, differential trace, and the integral identity.}
We give the comparison that computes the actual endomorphism
\(HF\), including the Prym components.  Fix an embedding
\(\overline\Q\hookrightarrow\C\).  For a finite curve
map \(q:C\to D\) of degree \(d\), the cotangent map
of \(q_*:J_C\to J_D\) is the usual pullback
\(q^*:H^0(D,\Omega^1)\to H^0(C,\Omega^1)\).
The cotangent map of divisor pullback
\(q^*:J_D\to J_C\) is the trace
\(q_!:H^0(C,\Omega^1)\to H^0(D,\Omega^1)\).
Indeed, on an unramified coordinate neighborhood, a moving
point of \(D\) pulls back to the sum of its \(d\)
moving inverse images.  Differentiating their Abel--Jacobi
integrals gives exactly the sum over the inverse branches,
which is \(q_!\).  The equality extends across the
finite branch locus.  The identification of the algebraic
Jacobian and these integrals is
\cite[Theorem 2.5 and its proof, pp.7--8]{RootMilneJV2021}.

On regular differentials use the positive Hermitian product
\[
       \langle\omega,\eta\rangle_C
          =\frac{i}{2}\int_{C(\C)}\omega\wedge\overline\eta.
\]
Integration over the inverse branches gives
\(\langle q^*\alpha,\omega\rangle_C
 =\langle\alpha,q_!\omega\rangle_D\), and
\(q_!q^*=[d]\).  Therefore the cotangent operator of
\(q^*q_*\) is \(d\) times the orthogonal projection
onto \(q^*H^0(D,\Omega^1)\).  This is a statement
about the specified trace and pullback, with both its
image and its multiplier determined.

Every automorphism \((u,v)\mapsto(\zeta u,\xi v)\)
with \(\zeta^8=\xi^8=1\) preserves this integral.
The forms \(\omega_{a,b}\) in (A.182.12) have
pairwise distinct characters \(\zeta^a\xi^b\).
For two distinct indices, choose an automorphism on
which their character ratio differs from one; invariance
of the product forces their inner product to be zero.
Thus these twenty-one character lines are mutually
orthogonal.

For each elliptic and genus-two factor the pullback
subspace is exactly its row of (A.182.13).  The operator
\(f_i^*f_{i*}\) consequently has cotangent eigenvalue
sixteen on its line and zero on all the other lines.
The operator \(c_jg_j^*g_{j*}\) has eigenvalue
\(c_j\deg g_j=16\) on its two lines and zero on
their orthogonal complement.

For a Prym factor, let
\(q_l:\operatorname{Jac}(D_l)\to\mathcal P_l\)
be \(1-\iota_l\) with its image as target, so that
\(j_lq_l=1-\iota_l\).  Its contribution to \(HF\)
is literally
\[
             h_l^*j_lq_lh_{l*}
                  =h_l^*(1-\iota_l)h_{l*}.
\tag{A.188.8}
\]
On differentials of \(D_l\), \(1-\iota_l^*\)
is twice the orthogonal projection onto the two
anti-invariant forms, since the involution preserves
the integral and has eigenvalues one and minus one.
The trace/pullback calculation with \(\deg h_l=8\)
makes (A.188.8) sixteen times the orthogonal projection
onto their pullbacks.  Those are precisely its two
lines in (A.182.13).  This computation uses the
inclusion into \(\operatorname{Jac}(D_l)\) and
requires no principal polarization on \(\mathcal P_l\).

The twelve displayed subspaces partition the full
twenty-one-line basis.  Summing their operators shows
that \(HF\) and \([16]\) have the same cotangent
map.  A homomorphism of complex abelian varieties is
determined by its differential: on the universal
covering vector spaces it is the induced complex
linear map.  Hence \(HF=[16]\) over \(\C\),
and therefore over \(\Q\), since both morphisms
were defined over \(\Q\).  Composing
\((FH-[16])F=0\) with the surjective isogeny \(F\)
proves \(FH=[16]\) as well.  The kernel assertions
now follow by applying these actual identities.  In
characteristic zero the degree of an isogeny is the
order of its finite kernel; \(J_8[16]\) is a
two-group, proving the degree assertion.

\emph{The explicit arithmetic bound.}
The maps in (A.188.1)--(A.188.2) induce the same
identities on continuous cocycles and on their
localization kernels.  A.187 proves
\(2\Sha(\mathcal B/\Q)[2^\infty]=0\).
For \(\xi\in\Sha(J_8/\Q)[2^\infty]\),
\[
                       32\xi=H(2F\xi)=0.
\]
The finite Kummer exact sequence (A.187.16) now
shows that the transition through five additional
levels kills its \(\Sha\)-coordinate.  Every image
is consequently a Mordell--Weil Kummer class.
Conversely, each such finite class is the image
of the same rational point at level \(n+5\).
This proves (A.188.3) with the asserted map and
every localization comparison.

The generic inverse equations, their forward
substitutions, the degrees and the twenty-one
operator multipliers are checked in
\path{code/point3_fermat8_retraction_certificate.py}.
The integration and trace argument above proves
the endomorphism identity itself.
\end{proof}
